\subsection{Thiết kế Kafka cho kiến trúc hướng sự kiện}

Trong hệ thống thương mại điện tử này, Kafka đóng vai trò là xương sống giao tiếp bất đồng bộ giữa các dịch vụ thành viên. Mô hình này giúp các dịch vụ hoạt động độc lập, giảm sự phụ thuộc trực tiếp qua giao thức HTTP, đồng thời hỗ trợ khả năng mở rộng, tính chống chịu lỗi và đảm bảo nhất quán dữ liệu giữa luồng ghi và luồng đọc trong kiến trúc CQRS.

\subsubsection{Vai trò của Kafka trong kiến trúc}

Kafka đảm nhận ba chức năng chính trong hệ thống:

\textbf{a/ Lớp luồng sự kiện (Event Streaming Layer)}

Các dịch vụ thành viên không gọi trực tiếp sang nhau khi xử lý các nghiệp vụ phức tạp, trừ trường hợp điều phối giao dịch. Quy trình vận hành như sau:
\begin{itemize}
    \item Dịch vụ nguồn phát ra sự kiện nghiệp vụ ngay sau khi cập nhật trạng thái thành công vào cơ sở dữ liệu chính.
    \item Các dịch vụ liên quan lắng nghe các chủ đề tương ứng để tiếp nhận dữ liệu.
    \item Mỗi dịch vụ xử lý phần việc nội bộ như gửi thông báo hoặc cập nhật kho dữ liệu AI mà không làm ảnh hưởng đến thời gian phản hồi của giao diện lập trình ứng dụng chính.
\end{itemize}

Ngoại lệ duy nhất là giao tiếp đồng bộ qua giao diện lập trình nội bộ được điều phối bởi Orchestrator trong các giao dịch phân tán Saga để đảm bảo tính nhất quán tức thời cho các bước quan trọng.

\textbf{b/ Đồng bộ hóa dữ liệu giữa các dịch vụ}

Trong mô hình CQRS, Kafka đóng vai trò then chốt trong việc lan truyền các thay đổi dữ liệu từ dịch vụ quản lý gốc sang các dịch vụ tiêu thụ khác. Ví dụ, khi người dùng cập nhật hồ sơ, dịch vụ AI Agent sẽ tiêu thụ sự kiện này để cập nhật lại thông tin trong cơ sở dữ liệu đồ thị Neo4j và cơ sở dữ liệu vector Milvus, đảm bảo dữ liệu phục vụ truy vấn luôn ở trạng thái mới nhất.

\textbf{c/ Lớp kết nối giữa các dịch vụ}

Thay vì tạo ra các kết nối chằng chịt giữa các thành phần, Kafka đóng vai trò là kênh sự kiện trung gian. Việc sử dụng mô hình công bố và đăng ký sự kiện giúp giảm thiểu sự phụ thuộc lẫn nhau, cho phép hệ thống dễ dàng tích hợp thêm các dịch vụ mới mà không cần thay đổi logic của các dịch vụ hiện có.

\subsubsection{Các chủ đề Kafka tiêu biểu trong hệ thống}

Để hỗ trợ kiến trúc hướng sự kiện, hệ thống sử dụng các nhóm chủ đề Kafka tiêu biểu nhằm truyền tải thay đổi trạng thái và dữ liệu nghiệp vụ:

\textbf{a/ Nhóm sự kiện người dùng}
\begin{itemize}
    \item \texttt{user.register.local}: Phát sinh khi người dùng tạo tài khoản mới.
    \item \texttt{user.login.local}: Ghi nhận sự kiện đăng nhập thành công.
    \item \texttt{user.updated}: Lan truyền thông tin cập nhật người dùng.
    \item \texttt{user.behavior.tracked}: Lưu vết hành vi phục vụ các tính năng thông minh.
\end{itemize}

\textbf{b/ Nhóm sự kiện giỏ hàng}
\begin{itemize}
    \item \texttt{cart.checked\_out}: Người dùng xác nhận thanh toán để Order Service tạo đơn hàng.
    \item \texttt{cart.updated}: Ghi nhận thay đổi giỏ hàng phục vụ phân tích sở thích.
\end{itemize}

\textbf{c/ Nhóm sự kiện đơn hàng}
\begin{itemize}
    \item \texttt{order.created}: Đơn hàng mới được tạo thành công.
    \item \texttt{order.status.updated}: Cập nhật trạng thái xử lý đơn hàng.
    \item \texttt{order.completed}: Đơn hàng hoàn tất để phục vụ thống kê báo cáo.
    \item \texttt{order.saga.failed}: Ghi nhận khi một giao dịch Saga bị hủy để kích hoạt luồng hoàn tác.
\end{itemize}

\textbf{d/ Nhóm sự kiện sản phẩm}
\begin{itemize}
    \item \texttt{product.created}: Sản phẩm mới được thêm vào hệ thống.
    \item \texttt{product.updated}: Thông tin sản phẩm được chỉnh sửa.
    \item \texttt{product.sales.updated}: Cập nhật dữ liệu bán hàng sau mỗi đơn hàng thành công.
\end{itemize}

\textbf{e/ Nhóm sự kiện quản trị}
\begin{itemize}
    \item \texttt{admin.log.created}: Lưu lại nhật ký hành động của quản trị viên.
    \item \texttt{admin.report.requested}: Tiếp nhận yêu cầu tạo báo cáo thống kê.
    \item \texttt{admin.report.generated}: Thông báo báo cáo đã hoàn thành và sẵn sàng truy xuất.
\end{itemize}

\textbf{f/ Nhóm sự kiện trí tuệ nhân tạo}
\begin{itemize}
    \item \texttt{ai.recommendation.requested}: Yêu cầu tạo gợi ý sản phẩm theo hồ sơ người dùng.
    \item \texttt{ai.recommendation.generated}: Kết quả gợi ý đã sẵn sàng trả về cho người dùng.
    \item \texttt{ai.user\_behavior.received}: AI tiếp nhận dữ liệu hành vi để huấn luyện mô hình.
\end{itemize}